% Section 7: the theory half of the draft's section 8 (notes/PLAN-article-structure.md, "The
% outlines"). The computation half is the computation paper's; nothing here cites it.
\section{Discussion}
\label{sec:discussion}

A pyramid applies one small mask at every level, and the classical criteria judge that mask on its
own \cite[p.~533]{burt1983laplacian}, \cite{meer1987frequency}. This paper judged it by another
standard. Covariance under one ratio $q$ requires the mask to be the step between two levels of a family
of kernels continuous in scale. With the input read as the image at the scale $1$, the levels of the
pyramid are then the image smoothed by the steps of that family from the scale $1$. The questions are which masks have such a family, and how much of it the mask fixes.

The new standard rules out every finite mask. A law of bounded support other than a point mass is
the step of no family: the cascade and continuity in scale exclude it
(Proposition~\ref{prop:no-bounded-step}). In its place
comes a criterion on the weights. A symmetric probability mask other than a point mass is the step of an exact
grid pyramid if and only if it is infinitely divisible with a logarithmic moment
(Proposition~\ref{prop:exact-masks}). For laws the sufficiency is the theorem of Maejima and Ueda
\cite[Thm.~3.3, p.~353]{maejima2009stochastic}, and the necessity of infinite divisibility comes from
continuity in scale. On a lattice such a mask has unbounded support
(Proposition~\ref{prop:exact-masks}(1)). It also has a positive transform, so at an integer ratio
it never has Burt's equal contribution, which requires zeros of the transform
(Section~\ref{subsec:grid-pyramids-line}). A pyramid that truncates its mask is therefore not
exact, and an
exact one must apply its mask without truncation. On the line a recursion does this. The sampled
Laplace mask is applied by one forward and one backward recursion
(Proposition~\ref{prop:sampled-laplace-exact}(2)). In every dimension the resolvent masks of the
discrete Helmholtz operators are exact as well (Proposition~\ref{prop:resolvent-masks-exact}).

The mask of an exact pyramid fixes the kernels at its levels. Label the scales so that the
dilation by $q$ maps the scale $t$ to $qt$. The kernel at a level is then the limit of the step
convolved with its dilates toward finer scales (Proposition~\ref{prop:ladder-knot-kernel}). This
is the equivalent weighting function of Burt and Adelson run downward in scale. For these steps its
limit exists and is the kernel. Between the levels a symmetric family is determined by its ladder,
up to a relabelling, exactly when the data of the step, its Gaussian part and \Levy{} measure, span
an extreme ray of the cone of such data (Proposition~\ref{prop:ladder-knot-kernel}(2)).

The strata place the pyramids among the families of scale space. The dilations under which a
symmetric family is covariant form $\{1\}$, some $q^{\ZZ}$, or all of $(0,\infty)$
(Lemma~\ref{lem:family-dilations}(2)). So covariance under the powers of one ratio is the only case
between none and every dilation. For isotropic families each stratum is strictly larger than the
next (Theorem~\ref{thm:lattice-strata-strict}). The pyramids of stratum~2 therefore form a wider
class than the fully covariant families that the isotropic paper classifies
\prog{iso v0.1}{Thm.~3.6, p.~13}. Lindeberg's scale space for discrete signals
\cite{lindeberg1990scalespace} lies in stratum~1 and not in stratum~2
(Proposition~\ref{prop:lindeberg-stratum-one} of the supplement). It is no pyramid in this sense.

Stratum~2 holds many families, so a designer still needs a reason to choose one. On the line the
count of local extrema selects among the fully covariant families \prog{V-selection v0.1}{\S4}.
In stratum~2 the rung-to-rung count holds on the ladders of a class
(Proposition~\ref{prop:ladder-non-creation-line}) and selects none of them. In the plane a bound on
the value at a local maximum takes its place. Among the isotropic families covariant under $q$, the
\Matern{} pyramid is, up to the gauge, the one whose every octave stage has the implicit-Euler
maximum principle (Theorem~\ref{thm:selection-stratum-two}(2)). The family selected is covariant
under every dilation, although only $q$ was required. Among symmetric rotation-invariant laws the
principle holds exactly for the Bessel potentials (Proposition~\ref{prop:stage-lemma}(3)), so it
rejects the Gaussian stages that classical non-enhancement selects on the line
\prog{V-selection v0.1}{Thm.~4.2(3), p.~31}. It is required of the octave stages only. A stage of
the \Matern{} pyramid has it only across a factor $q$ (Theorem~\ref{thm:selection-stratum-two}(3)),
and no symmetric family has it at every stage
(Proposition~\ref{prop:principle-other-readings}(2)). So each ratio $q$ selects its own fully
covariant family. Its octave step is one Helmholtz
resolvent, the law $L_{ct}$ (Theorem~\ref{thm:matern-pyramid}). The difference of the levels $t$
and $qt$ is $-c^2t^2$ times the Laplacian of the coarser one
(Proposition~\ref{prop:laplacian-pyramid}). So for this family the level difference of the
Laplacian pyramid of Burt and Adelson \cite[p.~535]{burt1983laplacian}, taken on the continuum before any
decimation, is a Laplacian.

Sequences of implicit heat steps with geometrically growing step sizes occur in nonstationary
iterated Tikhonov regularization \cite[pp.~38--39]{hanke1998nonstationary}. The decimation between
them was not found in that literature. Nor was a covariant family through them, unique up to the
gauge under the implicit-Euler principle. The line does not decide which isotropic families exist
in $\RR^d$ either. An isotropic family is fixed by its marginal (Lemma~\ref{lem:reduction-keeps-stratum}), but the
isotropic classes of one ratio shrink strictly with the dimension
(Proposition~\ref{prop:what-does-not-lift-chain}). Which families of the line extend to $\RR^d$
depends on $d$.

Computing on a lattice has a price. In $\RR^d$, $d \ge 2$, no kernel at a positive scale of an
exact grid pyramid whose mask is carried by a lattice is rotation invariant
(Proposition~\ref{prop:grid-pyramid-not-isotropic}(2)). Exact here is sense (iii) of
Definition~\ref{def:exact-senses}. The frame \Matern{} pyramid on the grid is exact and not
isotropic (Proposition~\ref{prop:grid-pyramid-not-isotropic}(3)). The \Matern{} pyramid is
isotropic, and its octave steps $L_{ct}$ have densities, so it carries the ladder of no mask on a
lattice. The obstruction rests on continuity in scale as much as on the lattice. Some symmetric
probability mask of infinite support on $\ZZ^d$, whose transform vanishes outside a small ball in
each period, has a rotation-invariant solution of the knot relation
(Proposition~\ref{prop:band-limited-refinable-mask}). Its dilates, applied at the knots, make an
isotropic two-parameter system of laws that is discontinuous at the knots. The mask is not
infinitely divisible, so no family continuous in scale carries it. The steerable pyramid uses
strictly band-limited low-pass filters, to prevent aliasing \cite[p.~445]{simoncelli1995steerable}.

\subsection{What the results do not say}

\emph{The families are positive and continuous in scale.} Every kernel here is a
probability law, so a mask with a negative weight is outside every statement. Burt's five-tap
generating kernel with central weight $a > 1/2$, which has negative outer weights, is one
\cite[pp.~533--534]{burt1983laplacian}. Symmetry is required where a result says so. In
Section~\ref{ch:price} it is required of the mask and not of the family that carries its ladder
(Definition~\ref{def:exact-grid-pyramid}), and the statements there are made for an arbitrary such
family.
Rotation invariance is required where a result says isotropic.
Proposition~\ref{prop:no-bounded-step} requires neither. The cascade together with continuity in
scale excludes the finite masks (Proposition~\ref{prop:no-bounded-step}). Without continuity the knot relation $D_q\mu_1 = \mu_1 * w$ alone
has solutions for a finite mask, such as the cubic B-spline for the binomial mask
(Proposition~\ref{prop:binomial-knot-relation}). In transforms the knot relation reads
$\widehat{\mu_1}(q\omega) = \widehat{w}(\omega)\widehat{\mu_1}(\omega)$. At $q = 2$ this is the
refinement equation of stationary subdivision, there with a mask of sum $2^d$
\cite[Ch.~6, p.~107]{cavaretta1991stationary}. The same equation defines the scaling functions of
wavelet theory \cite[p.~5]{cavaretta1991stationary}.

\emph{The selection does not count extrema.} In the plane the number of local maxima of a record is not
monotone in scale for any candidate treated here: it can rise at every rung of the \Matern{} and
frame \Matern{} pyramids and of the Gaussian ladder, for an input that depends on the rung
(Proposition~\ref{prop:creation-plane}).
Theorem~\ref{thm:selection-stratum-two} selects by a bound on the value of the output at a local
maximum. The counting statements that select on the line \prog{V-selection v0.1}{\S4} are not
offered as properties of a pyramid in the plane.

\emph{Exact means sense (iii).} ``Exact'' in Proposition~\ref{prop:grid-pyramid-not-isotropic} is
sense (iii) of Definition~\ref{def:exact-senses}: a family continuous in scale and covariant under
$q$ carries the ladder.
Whether a subsample of a grid family is again a member on the coarser grid, sense (ii), is a
different question. The paper answers it for one family only, on the line: the subsample of a
sampled Laplace kernel is proportional to the sampled Laplace kernel of the same range on the
coarser grid (Proposition~\ref{prop:sampled-laplace-exact}(2)). The range is that of the continuum
kernel sampled. The parameter of the discrete Helmholtz kernel with the same samples depends on
the spacing $h$. In $d \ge 2$ it is left to further
work.

\emph{The ratio is a scalar.} Outside
Proposition~\ref{prop:expanding-grid-pyramid-not-isotropic}, every dilation in the theory is
$x \mapsto qx$ on a lattice, the hexagonal lattice included. Some pyramids dilate by a map that also rotates the lattice: Crowley's
resampling by $\sqrt2$ \cite[\S 3.3.2, p.~32]{crowley1981representation}, the quincunx filter banks,
whose dilation matrix squares to $2I$ \cite[(2), p.~534]{kovacevic1992nonseparable}, and the
hexagonal pyramid of seven-point cells \cite[pp.~97, 99]{watson1989hexagonal}. For them only the
obstruction to isotropy is carried over
(Proposition~\ref{prop:expanding-grid-pyramid-not-isotropic}). Their families, strata and selection
are not treated.

\emph{The cost of computing in the plane is not treated.} The cost of a level in the plane, and the
distance of a decimated computation there from the continuum, are left to further work. On the
line the sampled Laplace family is an exact cascade on the grid, computed by one forward and one
backward first-order recursion (Propositions~\ref{prop:recursion-helmholtz}
and~\ref{prop:sampled-laplace-exact}). For an integer $q \ge 2$, the \Matern{} pyramid computed by
it with decimation agrees with the samples of the continuum steps at the first level and at no
later one, for every spacing (Proposition~\ref{prop:sampled-laplace-exact}(2)). The discrete
Helmholtz family on $h\ZZ^d$ is an exact cascade on the grid as well
(Proposition~\ref{prop:discrete-helmholtz}), but no recursion that computes it in the plane is
given here.
% TODO(module computation): cite the computation paper when released.

\subsection{Four questions the results answer}

Each main result leaves a question next to it: whether the creation of extrema can be placed at a
chosen rung, whether the principle survives decimation, whether a rotating dilation escapes the
obstruction, and whether a finite mask can keep isotropy by giving up exactness. Each has a short
answer from a statement of the paper.
\begin{enumerate}
\item \emph{Can the Gaussian ladder create extrema at a chosen rung?} Yes. For every $q > 1$ and
every $\tau > 0$, some positive input has a Gaussian record with one local maximum at the variance
$\tau$ and at least two at $q^2\tau$ (Proposition~\ref{prop:creation-plane}(5)).
\item \emph{Does the implicit-Euler principle survive decimation?} Read pointwise, no. The principle
holds for every octave stage of the \Matern{} pyramid on the continuum
(Theorem~\ref{thm:selection-stratum-two}(2)). On $h\ZZ^d$ the $(2d+1)$-point resolvent has it on
the axis neighbours (Proposition~\ref{prop:stage-lemma}(2)), and in $d \ge 2$ the frame mask has it
on no finite symmetric neighbourhood (Proposition~\ref{prop:grid-frame-step-not-selected}). After
decimation by an integer $q \ge 2$, one finitely supported input gives a decimated output with a
strict maximum over its axis neighbours on $qh\ZZ^d$, at a point where it exceeds the input
(Proposition~\ref{prop:principle-after-decimation}). There the inequality of the principle fails.
The principle concerns one smoothing step on one grid. Across a decimation the two arrays compared
are related by no convolution, since the decimated output depends on the input at the discarded
points. For the $(2d+1)$-point resolvent, at a point of the coarse grid that is also a maximum
over its axis neighbours on the fine grid, the inequality still holds, by
Proposition~\ref{prop:stage-lemma}(2).
\item \emph{Can a dilation that rotates the lattice give an isotropic exact grid pyramid?} No. In
$d \ge 2$ no exact grid pyramid whose dilation is an invertible linear map $M$ is isotropic
(Proposition~\ref{prop:expanding-grid-pyramid-not-isotropic}). This covers the similarities
$x \mapsto cRx$, $c > 1$ and $R \in \Orth(d)$, of the quincunx and hexagonal pyramids.
\item \emph{Can a finite mask have a rotation-invariant refinable law?} No. In $d \ge 2$ a
rotation-invariant law that satisfies the knot relation with a probability mask that has an
exponential moment, as every finite mask has, is $\delta_0$, and so is the mask. The same holds for a
finitely supported complex measure in place of the mask
(Proposition~\ref{prop:no-isotropic-refinable-finite-mask}). A mask of infinite support can: the
band-limited mask of Proposition~\ref{prop:band-limited-refinable-mask} has a rotation-invariant
refinable law, and its grid pyramid is not exact.
\end{enumerate}
